\section{A budget-matched evaluation protocol}\label{sec:C23:experiments}
The evaluation unit is a task plus initial environment, permitted resources, and scoring contract. Pin repository commit, dependency state, tool versions, model identity where available, and product configuration. If a hosted product does not expose a stable model or harness revision, record the retrieval and execution dates and restrict conclusions to the observed service. A nominally repeated experiment can change when the provider updates an opaque component.

Use task-level paired comparisons. Each system receives the same problem and a reset environment, with task order randomized where persistent memory is enabled. Declare whether memory may transfer between repositories or projects. A system with privileged workspace history has a different information budget from a fresh baseline. This difference may be a legitimate deployment advantage, but it must not be misrepresented as an intrinsic model gain.

Resource accounting includes input and output tokens, tool fees, environment startup, waiting, retries, verifier work, and human intervention. Report elapsed time and monetary cost separately. A cloud product can exploit parallel resources unavailable to a local baseline; matching elapsed time alone then compares different total compute. Conversely, matching dollars may permit different wall-clock and environment availability. The comparison axes should precede any ranking.

Ablations should target interfaces as well as prompts. Compare full repository search with restricted retrieval, executable feedback with model-only critique, persistent workspace with fresh workspace, and independent samples with coordinated subagents. Preserve test access and scoring across conditions. Where a product does not permit a component ablation, report the complete configured treatment rather than inventing an internal causal explanation.

Failure analysis should retain invalid calls, unexecuted tests, timeouts, partial patches, and inaccessible tasks. Distinguish infrastructure failures from policy errors without removing either from overall operational reliability. Confidence intervals should cluster repeated attempts by task and repository where applicable. If the sample is too small to distinguish adjacent products, report that uncertainty rather than converting a noisy ordering into a purchasing recommendation.

